\documentclass[]{scrartcl}

\usepackage{hyperref}
\usepackage{acro}
\pagenumbering{gobble}

\DeclareAcronym{sba}{
	short=SBA,
	long=Single Best Answer}
	
\DeclareAcronym{llm}{
	short=LLM,
	long=Large Language Model}

\begin{document}
\vspace*{-2.5cm}
\section{Introduction}
\acf{sba} questions are the dominant constituent of written postgraduate medical examinations, favoured for their reliability, ease of marking, and ability to test clinical reasoning.
Traditionally, their production has required considerable time and effort from clinicians with knowledge of the relevant context, which leads to scarcity, question reuse, and limited revision materials for candidates.

Developments in \acfp{llm} make them suitable for generation of text derived from standard sources of knowledge. Since their public release many authors have applied them to attempting examinations, however, no published literature exists applying them to the task of writing questions.

\section{Methods}
\textit{Minerva} is a command-line utility created for this project written in python, backed a cloud-based \ac{llm} interface, which makes use of a combination of retrieval-augmented generation, prompt engineering, few-short learning, and output structuring to produce \acp{sba} consistent with the target specification, such as the primary FRCA MCQ.

A quiz of 10 standard questions released by the RCoA, and 20 \ac{llm}-derived topic-matched questions created using the \textit{minerva} tool, was presented to anaesthetists, asking them to answer the question and assess the clarity, relevance, and difficulty on 5-point Likert scales, and to identify the source of the question.
Scores were tested for equivalence using a TOST procedure with a 0.5 equivalence bound.

\section{Results}
Responses were received from 8 anaesthetists with experience ranging from 6 months to over 10 years, half had passed the primary FRCA MCQ, for a total of 240 question evaluations.

The overall success rate was 63.8\% for the standard questions and 56.9\% for \ac{llm}-derived questions.
All metrics demonstrated statistically significant equivalence; clarity at 4.26$\pm$0.79 and 4.06$\pm$0.88 ($p < 0.01$), relevancy at 4.18$\pm$0.73 and 4.14$\pm$0.74 ($p < 0.0001$), and difficulty at  3.39$\pm$1.10 and 3.38$\pm$0.94 ($p < 0.001$) for standard and \ac{llm} questions respectively.
Standard questions were correctly identified in 50\% of cases (40/80), and respondents were not sure in 30\% (24/80), and in 20\% believed them to be \ac{llm}-derived. \ac{llm}-derived questions were correctly identified in 26\% of cases (42/160), not sure in 35\% (56/160), and in 39\% (62/160) believed the \ac{llm} questions were human written.

\section{Conclusion}
These results suggest that \ac{llm} techniques are capable of producing \acp{sba} that show a striking degree of similarity to reference questions released by the RCoA, as judged by anaesthetists. However, further work is needed to elicit whether they would be suitable for use in formalised examinations.
\newline
Source - \href{https://www.github.com/glfharris/minerva}{github.com/glfharris/minerva} under GPL-3.0.

\end{document}
